\section{Математическая модель рекомендательной системы с обратной связью}
\label{sec:Chapter2} \index{Chapter2}

% -----------------------------------------------------------------------
% 2.1 ЛОКАЛЬНАЯ ЛИНЕАРИЗАЦИЯ РЕКОМЕНДАТЕЛЬНОЙ ПОЛИТИКИ
% -----------------------------------------------------------------------

\subsection{Локальная линеаризация рекомендательной политики}
\label{subsec:local_linearization}

Пусть $\bar{w}$~--- параметры модели в стационарной точке и пусть
$f(u, i; w) \in \mathbb{R}$~--- скоринговая функция (оценка релевантности
объекта $i$ пользователю $u$ при параметрах $w$). Разложим $f$ по
параметрам вблизи $\bar{w}$:
\[
    f(u, i; w) \;\approx\; f(u, i; \bar{w}) \;+\;
    \bigl\langle \nabla_w f(u, i; \bar{w}),\, w - \bar{w} \bigr\rangle.
\]

Обозначим через $J_{ui} = \nabla_w f(u, i; \bar{w}) \in \mathbb{R}^p$
вектор-градиент скоринговой функции по параметрам в стационарной точке.
Тогда линеаризованная рекомендательная политика задаётся матрицей Якоби
$\mathbf{J} \in \mathbb{R}^{NM \times p}$, строки которой~--- это
$J_{ui}$ для всех пар $(u, i)$.

\begin{definition}[Локальный Якобиан рекомендательной политики]
\label{def:local_jacobian}
\textit{Локальным Якобианом} политики в кластере $k$ называется матрица
\[
    \mathbf{J}_k \;=\;
    \frac{1}{|\mathcal{C}_k||\mathcal{I}|}
    \sum_{u \in \mathcal{C}_k} \sum_{i \in \mathcal{I}}
    J_{ui} J_{ui}^\top \;\in\; \mathbb{R}^{p \times p},
\]
то есть усреднённая матрица внешних произведений градиентов по парам
пользователь-объект внутри кластера.
\end{definition}

Ключевое наблюдение: $\mathbf{J}_k$ отвечает за то, насколько
чувствительна рекомендательная политика к индивидуальным различиям
пользователей внутри кластера $k$. Если $\operatorname{rank}(\mathbf{J}_k)$
мало или $\mathbf{J}_k$ имеет маленький спектральный радиус,
политика практически не различает пользователей внутри кластера.

% -----------------------------------------------------------------------
% 2.2 КОВАРИАЦИОННАЯ ДИНАМИКА ВЫДАЧИ
% -----------------------------------------------------------------------

\subsection{Ковариационная динамика выдачи}
\label{subsec:covariance_dynamics}

Рассмотрим, как эволюционирует ковариационная матрица
$\Sigma_k^{(t)}$ под действием одного шага переобучения.

Пусть $\delta w^{(t)} = w^{(t+1)} - w^{(t)}$~--- изменение параметров
модели на шаге $t$. В линейном приближении это изменение пропорционально
стохастическому градиенту функционала потерь по наблюдаемым данным.
Для выдачи пользователя $u$ получаем:
\[
    \delta\hat{r}^{(t)}_u \;=\; \mathbf{J}^\top_u \,\delta w^{(t)},
\]
где $\mathbf{J}_u \in \mathbb{R}^{p \times M}$~--- матрица градиентов
по объектам для пользователя $u$.

\begin{proposition}[Рекуррентность ковариаций]
\label{prop:covariance_recursion}
В линейном приближении ковариационная матрица выдачи в кластере $k$
удовлетворяет рекуррентному соотношению:
\[
    \Sigma_k^{(t+1)} \;=\;
    \mathbf{A}_k\, \Sigma_k^{(t)}\, \mathbf{A}_k^\top
    \;+\; \mathbf{B}_k,
\]
где $\mathbf{A}_k$~--- \textup{эффективный оператор переобучения} для
кластера $k$, а $\mathbf{B}_k$~--- матрица, описывающая внешний шум
и exploration. Явные выражения для $\mathbf{A}_k$ и $\mathbf{B}_k$
выведены в приложении~\textup{A}.
\end{proposition}

Структура предложения~\ref{prop:covariance_recursion}~--- это матричное
уравнение дискретного Ляпунова. Долгосрочное поведение $\Sigma_k^{(t)}$
определяется спектральным радиусом $\rho(\mathbf{A}_k)$: если
$\rho(\mathbf{A}_k) < 1$, ковариации затухают; если $\rho(\mathbf{A}_k)
\geq 1$ при $\mathbf{B}_k = 0$~--- растут или стабилизируются.

% -----------------------------------------------------------------------
% 2.3 ЭФФЕКТИВНАЯ ДИНАМИКА ПЕРЕОБУЧЕНИЯ
% -----------------------------------------------------------------------

\subsection{Эффективная динамика переобучения}
\label{subsec:effective_retraining}

Связь между локальным Якобианом $\mathbf{J}_k$ и эффективным
оператором $\mathbf{A}_k$ устанавливается через матрицу Хессе
функционала потерь $\mathbf{H}^{(t)}$. При обучении методом градиентного
спуска с шагом $\eta$ получаем приближение:
\[
    \mathbf{A}_k \;\approx\; \mathbf{I} \;-\;
    \eta\, \mathbf{H}^{(t)}\big|_{\mathcal{C}_k}.
\]

Таким образом, спектральный радиус $\rho(\mathbf{A}_k)$ определяется
минимальными собственными значениями Хессе функционала потерь на
данных кластера $k$. Если Хессе вырожден или имеет малые собственные
значения (что типично для переобученных моделей), $\rho(\mathbf{A}_k)$
оказывается близким к $1$ или превышает $1$, и ковариации не затухают.

% -----------------------------------------------------------------------
% 2.4 ЗАМКНУТАЯ СИСТЕМА: ДИНАМИКА ПОЛЬЗОВАТЕЛЕЙ И EXPLORATION
% -----------------------------------------------------------------------

\subsection{Замкнутая система: динамика пользователей и exploration}
\label{subsec:full_system}

Чтобы модель была замкнутой и допускала анализ предельных режимов,
необходимо учесть два дополнительных элемента.

\textbf{Динамика пользователей.} Предполагается, что вектор предпочтений
пользователя $u$ медленно дрейфует под влиянием получаемых рекомендаций:
\[
    x_u^{(t+1)} \;=\; (1-\beta)\, x_u^{(t)} \;+\; \beta\,
    \operatorname{softmax}_{i}\!\bigl(f(u,i;w^{(t)})\bigr),
\]
где $\beta \in [0,1]$~--- коэффициент влияния рекомендаций на предпочтения.
При $\beta = 0$ пользователи статичны; при $\beta > 0$ их предпочтения
смещаются в сторону рекомендованных объектов.

\textbf{Параметр exploration.} В политику рекомендаций вводится
смешивание с равномерным распределением по объектам:
\[
    \pi^{(t)}_\alpha(i \mid u) \;=\;
    (1 - \alpha)\,\pi^{(t)}(i \mid u) \;+\; \frac{\alpha}{M},
    \quad \alpha \in [0,1].
\]
Параметр $\alpha$ контролирует степень исследования пространства объектов;
при $\alpha = 0$ система работает в чисто эксплуатационном режиме.

Полная замкнутая система описывается совместной эволюцией параметров
$w^{(t)}$, пользовательских предпочтений $\{x_u^{(t)}\}$ и политики
$\pi^{(t)}_\alpha$. Её анализ~--- в частности, условия сходимости к
стационарному распределению и существование предельных режимов с
ненулевым разнообразием~--- проводится в главе~\ref{sec:Chapter3}.
